% Abstract — two-phase framework with benchmark results (5-fold CV).
Brain tumor classification from magnetic resonance imaging (MRI) supports triage of glioma, meningioma, and pituitary lesions.
Most studies apply default parameters to handcrafted texture descriptors or couple exhaustive grid search to a single classifier, introducing optimisation bias and high compute cost.
We present a two-phase framework on a public MRI dataset of $3{,}064$ T1-weighted contrast-enhanced axial slices across three naturally imbalanced tumor classes (glioma, meningioma, pituitary): Phase~1 selects GLCM, LBP, DWT, and HOG hyperparameters using three classifier-free separability metrics (Fisher discriminant ratio, mutual information, Davies--Bouldin index) over 213 configurations; Phase~2 benchmarks eight classical models (SVM, KNN, random forest, XGBoost, LightGBM, logistic regression, MLP, ExtraTrees) with stratified 5-fold \texttt{GridSearchCV} on eight optimised feature compositions.
The best cross-validated configuration pairs a polynomial-kernel SVM with the fused $\mathrm{Full(opt)}$ representation (accuracy $=0.920$, macro-F1 $=0.920$, Cohen's $\kappa=0.88$); optimized HOG features alone reach macro-F1 $=0.903$ with SVM, and KNN on $\mathrm{Full(opt)}$ achieves macro-F1 $=0.900$.
Mean performance across classifiers confirms that fusion ($\mathrm{Full(opt)}$, mean F1 $=0.887$) outperforms individual descriptor families, with gradient descriptors ($D$, mean F1 $=0.865$) ranking highest among single-family sets.
Full statistical validation (Cohen's $\kappa$, bootstrap confidence intervals, McNemar, Wilcoxon, Friedman--Nemenyi) is reported for all 64 classifier--feature-set pairs.
These results show that separability-guided handcrafted features with rigorous classical benchmarking achieve strong MRI classification without deep learning or GPU training.
